\documentclass[final]{beamer}
% ====================
% Packages
% ====================

\usepackage[T1]{fontenc}
\usepackage{lmodern}
% change the orientation to landscape if needed
\usepackage[orientation=portrait,size=a0paper,scale=1.0]{beamerposter}
\usetheme{gemini}
\usecolortheme{uhb}
\usepackage{graphicx}
\usepackage{booktabs}
\usepackage{tikz}
\usepackage{pgfplots}
\pgfplotsset{compat=1.18}
\usepackage{anyfontsize}
\usepackage[style=alphabetic,backend=biber]{biblatex}
\addbibresource{poster.bib}
\usepackage{wrapfig}
\usepackage{subcaption}
\usepackage{setspace}
% ====================
% Lengths
% ====================
% 1.5 Zeilenabstand
%\onehalfspacing%

% If you have N columns, choose \sepwidth and \colwidth such that
% (N+1)*\sepwidth + N*\colwidth = \paperwidth
\newlength{\blocksep}
\newlength{\sepwidth}
\newlength{\colwidth}
\setlength{\blocksep}{0.05\pageheight}
\setlength{\sepwidth}{0.025\paperwidth}
\setlength{\colwidth}{0.3\paperwidth}

\newcommand{\separatorcolumn}{\begin{column}{\sepwidth}\end{column}}

% ====================
% Title
% ====================

\title{Land Use Change and Climate Change Impact on \newline Urban Heat Islands Development}

\author{Linus Andrae\inst{1}}

\institute[shortinst]{\inst{1} Institute for Environmental Physics, University of Bremen}

% ====================
% Footer (optional)
% ====================

\footercontent{%

  \begin{columns}[t]
    \begin{column}{\colwidth+5cm}
      \begin{wrapfigure}{l}{6cm}     
        \centering
        \vspace{-2cm}
        \includegraphics[width=5cm]{figures/qrcode.eps}
      \end{wrapfigure}
      \href{https://github.com/paradx/PresTPoster/blob/main/poster.pdf}{The online version of this Poster: \\ https://github.com/paradx/PresTPoster/blob/main/poster.pdf} \vspace{1cm}
    \end{column}
    \begin{column}{\colwidth}
      \vfill
      PresT Winter Semester 2023/2024 \\ University Bremen \hfill
    \end{column}
    \separatorcolumn%
    \begin{column}{\colwidth}
      \vfill
      \href{mailto:landrae@uni-bremen.de}{contact: landrae@uni-bremen.de}
    \end{column}
  \end{columns}
  }
% (can be left out to remove footer)

% ====================
% Logo (optional)
% ====================

% use this to include logos on the left and/or right side of the header:
\logoright{\includegraphics[height=5cm]{logos/logo}}
\logoleft{\includegraphics[height=5cm]{logos/logoIUP.png}}

% ===================
% Configuration of blocks
% ==================

\addtobeamertemplate{block begin}
  {}
  {\vspace{1ex plus 0.5ex minus 0.5ex} % Pads top of block
     % separates paragraphs in a block
   \setlength{\parskip}{24pt plus 1pt minus 1pt}%
   \begin{adjustwidth}{2cm}{2cm}
  }
\addtobeamertemplate{block end}
  {\end{adjustwidth}%
   \vspace{2ex plus 0.5ex minus 0.5ex}}% Pads bottom of block
  {\vspace{\blocksep}} % Seperates blocks from each other

\addtobeamertemplate{block alerted begin}
  {}
  {\vspace{1ex plus 0.5ex minus 0.5ex} % Pads top of block
     % separates paragraphs in a block
   \setlength{\parskip}{24pt plus 1pt minus 1pt}%
   \begin{adjustwidth}{2cm}{2cm}
  }
\addtobeamertemplate{block alerted end}
  {\end{adjustwidth}%
   \vspace{2ex plus 0.5ex minus 0.5ex}}% Pads bottom of block
  {\vspace{\blocksep}} % Seperates blocks from each other

\addtobeamertemplate{block example begin}
  {}
  {\vspace{1ex plus 0.5ex minus 0.5ex} % Pads top of block
     % separates paragraphs in a block
   \setlength{\parskip}{24pt plus 1pt minus 1pt}%
   \begin{adjustwidth}{2cm}{2cm}
  }
\addtobeamertemplate{block example end}
  {\end{adjustwidth}%
   \vspace{2ex plus 0.5ex minus 0.5ex}}% Pads bottom of block
  {\vspace{\blocksep}} % Seperates blocks from each other

% ====================
% Body
% ====================

\begin{document}

\begin{frame}[t]
\vspace{1em}
\tikz[remember picture, overlay] \node[opacity=0.3, yshift=-15mm] at (current page.center){\includegraphics[width=\textwidth,height=0.99\textheight]{figures/overlay}};
\begin{columns}[t]
\separatorcolumn%

\begin{column}{\colwidth}

  \begin{block}{Abstract}
    Urban Heat Islands are a phenomena where the temperature within a city is significantly higher within an urban area, compared to the temperature of the surrounding rural area. \\
    Detection using remote sensing data is used frequently. 
    This Work adapted and combined different methods, to measure the impact of land use/ land cover change and longer term temperature changes on Urban Heat Islands in different European cities.
    
  \end{block}

  \begin{exampleblock}{Processing Chain}
   The Image Processing Pipeline can be separated in these steps:
    \begin{itemize}
      \item \textbf{Preprocessing} Cutting the image area to the Area of Interested
      \item \textbf{Feature Engineering} A Gabor filter bank of rotated Gabor kernels is used to enrich pixels with information about there place in larger structures of the image
      \item \textbf{Image Classification}  The pixels are classified using a random forest machine learning model that was trained using k-means based classification. 
      \item \textbf{Urban Area Detection} Classes indicating Urban Areas are identified and the buffer zones are build up around the areas to calculate a reference temperature.
      \item \textbf{Peri-Urban Area Clean-up} Larger structures of small settlements are masked away from the surrounding areas
      \item \textbf{Detection of Urban Heat Islands} Using the Thermal Infrared channels of Landsat, Urban heat islands are detected within the Urban Area, using a minimum threshold above the rural surrounding area.
      \item \textbf{Statistical Analysis of UHIs} Pixel within the UHI and around are evaluated using land surface type, NDVI and NDBI indicies
      \item \textbf{Timeline Creation} Repeating this process for different images of the same area over multiple years, change in extend and size and statistical composition of the UHIs can be analysed. This allows an inside in the impact of urbanisation and change in surface type on UHIs.
      \item \textbf{Climate Change Impact} Estimation of the imapact of rising average temperatures is done by substracting the statistically significant impact of metrological and surface change on the size of urban heat islands and correlating the UHI size with measurement from weather stations in the area in question.
    \end{itemize}
  \end{exampleblock}

  \begin{block}{Feature Engineering}
    
     \begin{figure}
       \begin{subfigure}{0.34\textwidth}
      \centering
      \includegraphics[width=\textwidth]{figures/gabor}
      \subcaption{Gabor Filter Bank~\cite{mediumgabor}} 
    \end{subfigure}
    \begin{subfigure}{0.34\textwidth}
      \centering
      \includegraphics[width=\textwidth]{figures/features}
      \subcaption{Features extracted from an image of the city of Bremen}
    \end{subfigure}
    \end{figure}

  \end{block}

\end{column}

\separatorcolumn%

\begin{column}{\colwidth}

  \begin{block}{A block containing an enumerated list}

    Vivamus congue volutpat elit non semper. Praesent molestie nec erat ac
    interdum. In quis suscipit erat. \textbf{Phasellus mauris felis, molestie
    ac pharetra quis}, tempus nec ante. Donec finibus ante vel purus mollis
    fermentum. Sed felis mi, pharetra eget nibh a, feugiat eleifend dolor. Nam
    mollis condimentum purus quis sodales. Nullam eu felis eu nulla eleifend
    bibendum nec eu lorem. Vivamus felis velit, volutpat ut facilisis ac,
    commodo in metus.

    \begin{enumerate}
      \item \textbf{Morbi mauris purus}, egestas at vehicula et, convallis
        accumsan orci. Orci varius natoque penatibus et magnis dis parturient
        montes, nascetur ridiculus mus.
      \item \textbf{Cras vehicula blandit urna ut maximus}. Aliquam blandit nec
        massa ac sollicitudin. Curabitur cursus, metus nec imperdiet bibendum,
        velit lectus faucibus dolor, quis gravida metus mauris gravida turpis.
      \item \textbf{Vestibulum et massa diam}. Phasellus fermentum augue non
        nulla accumsan, non rhoncus lectus condimentum.
    \end{enumerate}

  \end{block}

  \begin{block}{Fusce aliquam magna velit}

    Et rutrum ex euismod vel. Pellentesque ultricies, velit in fermentum
    vestibulum, lectus nisi pretium nibh, sit amet aliquam lectus augue vel
    velit. Suspendisse rhoncus massa porttitor augue feugiat molestie. Sed
    molestie ut orci nec malesuada. Sed ultricies feugiat est fringilla
    posuere.

    \begin{figure}
      \centering
      \begin{tikzpicture}
        \begin{axis}[
            scale only axis,
            no markers,
            domain=0:2*pi,
            samples=100,
            axis lines=center,
            axis line style={-},
            ticks=none]
          \addplot[red] {sin(deg(x))};
          \addplot[blue] {cos(deg(x))};
        \end{axis}
      \end{tikzpicture}
      \caption{Another figure caption.}
    \end{figure}

  \end{block}

  \begin{block}{Nam cursus consequat egestas}

    Nulla eget sem quam. Ut aliquam volutpat nisi vestibulum convallis. Nunc a
    lectus et eros facilisis hendrerit eu non urna. Interdum et malesuada fames
    ac ante \textit{ipsum primis} in faucibus. Etiam sit amet velit eget sem
    euismod tristique. Praesent enim erat, porta vel mattis sed, pharetra sed
    ipsum. Morbi commodo condimentum massa, \textit{tempus venenatis} massa
    hendrerit quis. Maecenas sed porta est. Praesent mollis interdum lectus,
    sit amet sollicitudin risus tincidunt non.

    Etiam sit amet tempus lorem, aliquet condimentum velit. Donec et nibh
    consequat, sagittis ex eget, dictum orci. Etiam quis semper ante. Ut eu
    mauris purus. Proin nec consectetur ligula. Mauris pretium molestie
    ullamcorper. Integer nisi neque, aliquet et odio non, sagittis porta justo.

    \begin{itemize}
      \item \textbf{Sed consequat} id ante vel efficitur. Praesent congue massa
        sed est scelerisque, elementum mollis augue iaculis.
        \begin{itemize}
          \item In sed est finibus, vulputate
            nunc gravida, pulvinar lorem. In maximus nunc dolor, sed auctor eros
            porttitor quis.
          \item Fusce ornare dignissim nisi. Nam sit amet risus vel lacus
            tempor tincidunt eu a arcu.
          \item Donec rhoncus vestibulum erat, quis aliquam leo
            gravida egestas.
        \end{itemize}
      \item \textbf{Sed luctus, elit sit amet} dictum maximus, diam dolor
        faucibus purus, sed lobortis justo erat id turpis.
      \item \textbf{Pellentesque facilisis dolor in leo} bibendum congue.
        Maecenas congue finibus justo, vitae eleifend urna facilisis at.
    \end{itemize}

  \end{block}

\end{column}

\separatorcolumn

\begin{column}{\colwidth}

  \begin{exampleblock}{A highlighted block containing some math}

    A different kind of highlighted block.

    $$
    \int_{-\infty}^{\infty} e^{-x^2}\,dx = \sqrt{\pi}
    $$

    Interdum et malesuada fames $\{1, 4, 9, \ldots\}$ ac ante ipsum primis in
    faucibus. Cras eleifend dolor eu nulla suscipit suscipit. Sed lobortis non
    felis id vulputate.

    \heading{A heading inside a block}

    Praesent consectetur mi $x^2 + y^2$ metus, nec vestibulum justo viverra
    nec. Proin eget nulla pretium, egestas magna aliquam, mollis neque. Vivamus
    dictum $\mathbf{u}^\intercal\mathbf{v}$ sagittis odio, vel porta erat
    congue sed. Maecenas ut dolor quis arcu auctor porttitor.

    \heading{Another heading inside a block}

    Sed augue erat, scelerisque a purus ultricies, placerat porttitor neque.
    Donec $P(y \mid x)$ fermentum consectetur $\nabla_x P(y \mid x)$ sapien
    sagittis egestas. Duis eget leo euismod nunc viverra imperdiet nec id
    justo.

  \end{exampleblock}

  \begin{block}{Nullam vel erat at velit convallis laoreet}

    Class aptent taciti sociosqu ad litora torquent per conubia nostra, per
    inceptos himenaeos. Phasellus libero enim, gravida sed erat sit amet,
    scelerisque congue diam. Fusce dapibus dui ut augue pulvinar iaculis.

    \begin{table}
      \centering
      \begin{tabular}{l r r c}
        \toprule
        \textbf{First column} & \textbf{Second column} & \textbf{Third column} & \textbf{Fourth} \\
        \midrule
        Foo & 13.37 & 384,394 & $\alpha$ \\
        Bar & 2.17 & 1,392 & $\beta$ \\
        Baz & 3.14 & 83,742 & $\delta$ \\
        Qux & 7.59 & 974 & $\gamma$ \\
        \bottomrule
      \end{tabular}
      \caption{A table caption.}
    \end{table}

    Donec quis posuere ligula. Nunc feugiat elit a mi malesuada consequat. Sed
    imperdiet augue ac nibh aliquet tristique. Aenean eu tortor vulputate,
    eleifend lorem in, dictum urna. Proin auctor ante in augue tincidunt
    tempor. Proin pellentesque vulputate odio, ac gravida nulla posuere
    efficitur. Aenean at velit vel dolor blandit molestie. Mauris laoreet
    commodo quam, non luctus nibh ullamcorper in. Class aptent taciti sociosqu
    ad litora torquent per conubia nostra, per inceptos himenaeos.

    Nulla varius finibus volutpat. Mauris molestie lorem tincidunt, iaculis
    libero at, gravida ante. Phasellus at felis eu neque suscipit suscipit.
    Integer ullamcorper, dui nec pretium ornare, urna dolor consequat libero,
    in feugiat elit lorem euismod lacus. Pellentesque sit amet dolor mollis,
    auctor urna non, tempus sem.

  \end{block}

  \begin{block}{References}

    \nocite{*}
    \footnotesize{\printbibliography}

  \end{block}

\end{column}

\separatorcolumn
\end{columns}
\end{frame}

\end{document}
